% =========================================================
% 3.10.6 Binary Adder
% =========================================================

\subsection{বাইনারি অ্যাডার
(\textenglish{Binary Adder})}

\textenglish{Half Adder} ও \textenglish{Full Adder} মূলত এক-বিট
বাইনারি সংখ্যা যোগ করার মৌলিক সার্কিট। কিন্তু কম্পিউটারে সাধারণত
\(4\)-\textenglish{bit}, \(8\)-\textenglish{bit},
\(16\)-\textenglish{bit}, \(32\)-\textenglish{bit} ইত্যাদি বহু-বিট
সংখ্যা নিয়ে কাজ করতে হয়।

যেমন,

\[
(1011)_2+(0110)_2
\]

যোগ করতে হলে শুধু এক জোড়া \textenglish{bit} যোগ করলেই হবে না।
প্রতিটি \textenglish{bit-position}-এর \textenglish{bit} যোগ করতে হবে
এবং এক ধাপে উৎপন্ন \textenglish{Carry} পরবর্তী ধাপে পাঠাতে হবে।

একাধিক \textenglish{bit}-এর দুটি বাইনারি সংখ্যা যোগ করার জন্য
একাধিক \textenglish{Adder} প্রয়োজনমতো সংযুক্ত করে যে সার্কিট তৈরি করা হয়,
তাকে \textbf{\textenglish{Binary Adder}} বলে।

\subsubsection*{সহজভাবে}

\[
\boxed{
\text{\textenglish{Multi-bit Binary Numbers}}
\rightarrow
\text{\textenglish{Binary Adder}}
\rightarrow
\text{\textenglish{Sum + Final Carry}}
}
\]


% =========================================================
% Types of Binary Adder
% =========================================================

\subsection{বাইনারি অ্যাডারের প্রকারভেদ}

বাইনারি সংখ্যার বিটগুলো কীভাবে অ্যাডারে প্রবেশ করে এবং কীভাবে যোগ
সম্পন্ন হয়, তার ভিত্তিতে \textenglish{Binary Adder} প্রধানত দুই প্রকার—

\begin{enumerate}[leftmargin=*,itemsep=2pt]
    \item \textbf{\textenglish{Serial Binary Adder}}
    \item \textbf{\textenglish{Parallel Binary Adder}}
\end{enumerate}


% =========================================================
% Serial Binary Adder
% =========================================================

\paragraph{\textenglish{Serial Binary Adder}:}

যে বাইনারি অ্যাডারে দুটি বহু-বিট বাইনারি সংখ্যার সংশ্লিষ্ট বিটগুলো
একটির পর একটি ধারাবাহিকভাবে যোগ করা হয়,
তাকে \textenglish{Serial Binary Adder} বলে।

এখানে সাধারণত একটি \textenglish{Full Adder} ব্যবহার করে
প্রতিটি ধাপে এক জোড়া বিট যোগ করা হয়।
এক ধাপে উৎপন্ন \textenglish{Carry} পরবর্তী ধাপের জন্য সংরক্ষণ করা হয়।

ধরা যাক,

\[
A=A_3A_2A_1A_0
\]

এবং,

\[
B=B_3B_2B_1B_0
\]

তাহলে পর্যায়ক্রমে যোগ হবে—

\[
A_0+B_0+C_0
\]

\[
A_1+B_1+C_1
\]

\[
A_2+B_2+C_2
\]

\[
A_3+B_3+C_3
\]

\par\begingroup

\textbf{মনে রাখবে:}
\textenglish{Serial Binary Adder}-এ একই \textenglish{Full Adder}
বারবার ব্যবহার করে একটির পর একটি বিট যোগ করা যায়।
তাই হার্ডওয়্যার কম লাগে, কিন্তু যোগ সম্পন্ন করতে তুলনামূলক বেশি সময় লাগে।
\par\endgroup



% =========================================================
% Parallel Binary Adder
% =========================================================

\paragraph{\textenglish{Parallel Binary Adder}:}

যে বাইনারি অ্যাডারে দুটি বহু-বিট বাইনারি সংখ্যার সংশ্লিষ্ট বিটগুলো
একই সময়ে পৃথক \textenglish{Adder stage}-এ প্রবেশ করে,
তাকে \textenglish{Parallel Binary Adder} বলে।

প্রতিটি \textenglish{bit-position}-এর জন্য আলাদা
\textenglish{Full Adder} ব্যবহার করা হয়।

যেমন, দুটি \(4\)-\textenglish{bit} বাইনারি সংখ্যা যোগ করতে
চারটি \textenglish{Full Adder} ধারাবাহিকভাবে সংযুক্ত করা যায়।


% =========================================================
% 4-bit Parallel Binary Adder
% =========================================================

\subsubsection*{৪-বিট প্যারালাল বাইনারি অ্যাডার
(\textenglish{4-bit Parallel Binary Adder})}

ধরা যাক, দুটি \(4\)-\textenglish{bit} বাইনারি সংখ্যা—

\[
A_3A_2A_1A_0
\]

এবং,

\[
B_3B_2B_1B_0
\]

যোগ করতে হবে।

এক্ষেত্রে চারটি \textenglish{Full Adder} ব্যবহার করা হয়।
সর্বডানের \textenglish{Full Adder} সর্বনিম্ন গুরুত্বপূর্ণ বিট
(\textenglish{LSB}) এবং সর্ববামের \textenglish{Full Adder}
সর্বাধিক গুরুত্বপূর্ণ বিট (\textenglish{MSB}) যোগ করে।

প্রাথমিকভাবে \textenglish{Carry-in} না থাকলে—

\[
C_0=0
\]

ধরা হয়।

প্রতিটি ধাপে—

\[
FA_0:
A_0+B_0+C_0
\longrightarrow
S_0,\;C_1
\]

\[
FA_1:
A_1+B_1+C_1
\longrightarrow
S_1,\;C_2
\]

\[
FA_2:
A_2+B_2+C_2
\longrightarrow
S_2,\;C_3
\]

\[
FA_3:
A_3+B_3+C_3
\longrightarrow
S_3,\;C_4
\]

অতএব সম্পূর্ণ ফলাফল—

\[
\boxed{C_4S_3S_2S_1S_0}
\]


% =========================================================
% 4-bit Parallel Binary Adder Diagram
% =========================================================

\begin{center}

\begin{chThreeDiagram}
\begin{tikzpicture}[
    >=Stealth,
    line width=0.8pt,
    fa/.style={
        draw,
        rounded corners=2pt,
        fill=ICTLightBlue,
        minimum width=2.25cm,
        minimum height=1.15cm,
        align=center
    }
]

\node[fa] (fa3) at (0,0) {\textbf{\(FA_3\)}};
\node[fa] (fa2) at (3.6,0) {\textbf{\(FA_2\)}};
\node[fa] (fa1) at (7.2,0) {\textbf{\(FA_1\)}};
\node[fa] (fa0) at (10.8,0) {\textbf{\(FA_0\)}};

% A3, B3
\draw[->] (-0.55,2.0) -- (-0.55,0.62);
\node[above] at (-0.55,2.0) {\(A_3\)};

\draw[->] (0.55,2.0) -- (0.55,0.62);
\node[above] at (0.55,2.0) {\(B_3\)};

% A2, B2
\draw[->] (3.05,2.0) -- (3.05,0.62);
\node[above] at (3.05,2.0) {\(A_2\)};

\draw[->] (4.15,2.0) -- (4.15,0.62);
\node[above] at (4.15,2.0) {\(B_2\)};

% A1, B1
\draw[->] (6.65,2.0) -- (6.65,0.62);
\node[above] at (6.65,2.0) {\(A_1\)};

\draw[->] (7.75,2.0) -- (7.75,0.62);
\node[above] at (7.75,2.0) {\(B_1\)};

% A0, B0
\draw[->] (10.25,2.0) -- (10.25,0.62);
\node[above] at (10.25,2.0) {\(A_0\)};

\draw[->] (11.35,2.0) -- (11.35,0.62);
\node[above] at (11.35,2.0) {\(B_0\)};

% Carry chain
\draw[<-] (fa3.west) -- ++(-1.25,0);
\node[above] at (-1.75,0.08) {\(C_4\)};

\draw[<-] (fa3.east) -- (fa2.west);
\node[above] at (1.8,0.08) {\(C_3\)};

\draw[<-] (fa2.east) -- (fa1.west);
\node[above] at (5.4,0.08) {\(C_2\)};

\draw[<-] (fa1.east) -- (fa0.west);
\node[above] at (9.0,0.08) {\(C_1\)};

\draw[<-] (fa0.east) -- ++(1.30,0);
\node[above] at (12.55,0.08) {\(C_0=0\)};

% Sum outputs
\draw[->] (fa3.south) -- ++(0,-1.15);
\node[below] at (0,-1.15) {\(S_3\)};

\draw[->] (fa2.south) -- ++(0,-1.15);
\node[below] at (3.6,-1.15) {\(S_2\)};

\draw[->] (fa1.south) -- ++(0,-1.15);
\node[below] at (7.2,-1.15) {\(S_1\)};

\draw[->] (fa0.south) -- ++(0,-1.15);
\node[below] at (10.8,-1.15) {\(S_0\)};

\end{tikzpicture}
\end{chThreeDiagram}

\end{center}

\[
\boxed{
\text{\textenglish{Output}}
=
C_4S_3S_2S_1S_0
}
\]

এখানে \(A_3A_2A_1A_0\) এবং \(B_3B_2B_1B_0\) হলো দুটি
\(4\)-\textenglish{bit} বাইনারি সংখ্যা।
\(S_3S_2S_1S_0\) হলো চারটি \textenglish{Sum bit}
এবং \(C_4\) হলো চূড়ান্ত \textenglish{Carry}।


% =========================================================
% Carry Propagation
% =========================================================

\subsubsection*{ক্যারি প্রবাহ
(\textenglish{Carry Propagation})}

প্রতিটি \textenglish{Full Adder}-এর \textenglish{Carry-out}
পরবর্তী \textenglish{Full Adder}-এর \textenglish{Carry-in}
হিসেবে প্রবেশ করে।

অর্থাৎ,

\[
C_0
\rightarrow
FA_0
\rightarrow
C_1
\rightarrow
FA_1
\rightarrow
C_2
\rightarrow
FA_2
\rightarrow
C_3
\rightarrow
FA_3
\rightarrow
C_4
\]

এইভাবে \textenglish{Carry} এক \textenglish{stage} থেকে
পরবর্তী \textenglish{stage}-এ ক্রমান্বয়ে ছড়িয়ে যায়।

এই কারণে এ ধরনের সাধারণ \textenglish{Parallel Adder}-কে
\textenglish{Ripple Carry Adder}-ও বলা হয়।

\par\begingroup

\textbf{মনে রাখবে:}
\textenglish{Parallel Adder}-এ \textenglish{input bit}গুলো একই সময়ে
সংশ্লিষ্ট \textenglish{Full Adder}-এ প্রবেশ করলেও
\textenglish{Carry} এক \textenglish{stage} থেকে পরবর্তী
\textenglish{stage}-এ ধারাবাহিকভাবে প্রবাহিত হয়।
\par\endgroup



% =========================================================
% Example: 1011 + 1111
% =========================================================

\subsubsection*{উদাহরণ:
\((1011)_2+(1111)_2\)}

ধরা যাক,

\[
A=(1011)_2
\]

এবং,

\[
B=(1111)_2
\]

তাহলে—

\[
A_3A_2A_1A_0=1011
\]

এবং,

\[
B_3B_2B_1B_0=1111
\]

প্রাথমিক \textenglish{Carry}—

\[
C_0=0
\]


% =========================================================
% Example Diagram
% =========================================================

\begin{center}

\begin{chThreeDiagram}
\begin{tikzpicture}[
    >=Stealth,
    line width=0.8pt,
    fa/.style={
        draw,
        rounded corners=2pt,
        fill=ICTLightGreen,
        minimum width=2.25cm,
        minimum height=1.15cm,
        align=center
    }
]

\node[fa] (fa3) at (0,0) {\textbf{\(FA_3\)}};
\node[fa] (fa2) at (3.6,0) {\textbf{\(FA_2\)}};
\node[fa] (fa1) at (7.2,0) {\textbf{\(FA_1\)}};
\node[fa] (fa0) at (10.8,0) {\textbf{\(FA_0\)}};

% FA3
\node[above] at (-0.55,2.0) {\(A_3=1\)};
\node[above] at (0.60,2.0) {\(B_3=1\)};
\draw[->] (-0.55,1.72) -- (-0.55,0.62);
\draw[->] (0.60,1.72) -- (0.60,0.62);

% FA2
\node[above] at (3.05,2.0) {\(A_2=0\)};
\node[above] at (4.20,2.0) {\(B_2=1\)};
\draw[->] (3.05,1.72) -- (3.05,0.62);
\draw[->] (4.20,1.72) -- (4.20,0.62);

% FA1
\node[above] at (6.65,2.0) {\(A_1=1\)};
\node[above] at (7.80,2.0) {\(B_1=1\)};
\draw[->] (6.65,1.72) -- (6.65,0.62);
\draw[->] (7.80,1.72) -- (7.80,0.62);

% FA0
\node[above] at (10.25,2.0) {\(A_0=1\)};
\node[above] at (11.40,2.0) {\(B_0=1\)};
\draw[->] (10.25,1.72) -- (10.25,0.62);
\draw[->] (11.40,1.72) -- (11.40,0.62);

% Carry values
\draw[<-] (fa3.west) -- ++(-1.20,0);
\node[above] at (-1.70,0.08) {\(C_4=1\)};

\draw[<-] (fa3.east) -- (fa2.west);
\node[above] at (1.80,0.08) {\(C_3=1\)};

\draw[<-] (fa2.east) -- (fa1.west);
\node[above] at (5.40,0.08) {\(C_2=1\)};

\draw[<-] (fa1.east) -- (fa0.west);
\node[above] at (9.00,0.08) {\(C_1=1\)};

\draw[<-] (fa0.east) -- ++(1.20,0);
\node[above] at (12.50,0.08) {\(C_0=0\)};

% Sum outputs
\draw[->] (fa3.south) -- ++(0,-1.15);
\node[below] at (0,-1.15) {\(S_3=1\)};

\draw[->] (fa2.south) -- ++(0,-1.15);
\node[below] at (3.6,-1.15) {\(S_2=0\)};

\draw[->] (fa1.south) -- ++(0,-1.15);
\node[below] at (7.2,-1.15) {\(S_1=1\)};

\draw[->] (fa0.south) -- ++(0,-1.15);
\node[below] at (10.8,-1.15) {\(S_0=0\)};

\end{tikzpicture}
\end{chThreeDiagram}

\end{center}


% =========================================================
% Stage-by-stage Calculation
% =========================================================

\par\begingroup


\begin{minipage}[t]{0.235\textwidth}
\centering

\textbf{স্টেজ–০}\\[-1mm]
{\footnotesize \textenglish{LSB}}

\vspace{2mm}

\[
A_0=1,\quad B_0=1
\]

\[
C_0=0
\]

\[
1+1+0=10
\]

\[
\boxed{S_0=0,\;C_1=1}
\]

\end{minipage}
\hfill
\begin{minipage}[t]{0.235\textwidth}
\centering

\textbf{স্টেজ–১}

\vspace{4mm}

\[
A_1=1,\quad B_1=1
\]

\[
C_1=1
\]

\[
1+1+1=11
\]

\[
\boxed{S_1=1,\;C_2=1}
\]

\end{minipage}
\hfill
\begin{minipage}[t]{0.235\textwidth}
\centering

\textbf{স্টেজ–২}

\vspace{4mm}

\[
A_2=0,\quad B_2=1
\]

\[
C_2=1
\]

\[
0+1+1=10
\]

\[
\boxed{S_2=0,\;C_3=1}
\]

\end{minipage}
\hfill
\begin{minipage}[t]{0.235\textwidth}
\centering

\textbf{স্টেজ–৩}\\[-1mm]
{\footnotesize \textenglish{MSB}}

\vspace{2mm}

\[
A_3=1,\quad B_3=1
\]

\[
C_3=1
\]

\[
1+1+1=11
\]

\[
\boxed{S_3=1,\;C_4=1}
\]

\end{minipage}

\par\endgroup

অতএব,

\[
C_4S_3S_2S_1S_0=11010
\]

সুতরাং,

\[
\boxed{
(1011)_2+(1111)_2=(11010)_2
}
\]

দশমিক পদ্ধতিতে যাচাই করলে—

\[
(1011)_2=(11)_{10}
\]

এবং,

\[
(1111)_2=(15)_{10}
\]

তাই,

\[
11+15=26
\]

এবং,

\[
(26)_{10}=(11010)_2
\]

অতএব ফল সঠিক।


% =========================================================
% Serial vs Parallel Comparison
% =========================================================

\subsubsection*{\textenglish{Serial} ও \textenglish{Parallel Binary Adder}-এর পার্থক্য}

\begin{center}

\renewcommand{\arraystretch}{1.25}

\begin{tabularx}
{\textwidth}
{|L{0.20\textwidth}|Y|Y|}

\hline

\rowcolor{ICTLightBlue}

\textbf{বিষয়}
&
\textbf{\textenglish{Serial Binary Adder}}
&
\textbf{\textenglish{Parallel Binary Adder}}
\\
\hline

যোগের পদ্ধতি
&
বিট একটির পর একটি যোগ করা হয়।
&
সংশ্লিষ্ট বিটগুলো একই সময়ে পৃথক
\textenglish{Adder}-এ প্রবেশ করে।
\\
\hline

\textenglish{Adder} সংখ্যা
&
তুলনামূলক কম।
&
তুলনামূলক বেশি।
\\
\hline

৪-বিট যোগে
&
সাধারণত একই \textenglish{Full Adder} বারবার ব্যবহার করা হয়।
&
সাধারণত ৪টি \textenglish{Full Adder} ব্যবহার করা হয়।
\\
\hline

সময়
&
বেশি লাগে।
&
তুলনামূলক কম লাগে।
\\
\hline

\textenglish{Hardware}
&
কম প্রয়োজন।
&
বেশি প্রয়োজন।
\\
\hline

গতি
&
কম।
&
বেশি।
\\
\hline

\end{tabularx}

\end{center}


% =========================================================
% Advantages and Limitations
% =========================================================

\subsubsection*{\textenglish{Parallel Binary Adder}-এর সুবিধা ও সীমাবদ্ধতা}

\textbf{সুবিধা:}

\begin{itemize}[label={{\bulletfont\textbullet}}, leftmargin=*, nosep]

    \item একাধিক বিট একই সময়ে সংশ্লিষ্ট
    \textenglish{Full Adder}-এ প্রবেশ করতে পারে।

    \item \textenglish{Serial Adder}-এর তুলনায় দ্রুত।

    \item বহু-বিট বাইনারি সংখ্যা সহজে যোগ করা যায়।

    \item \textenglish{Processor} ও
    \textenglish{ALU}-তে ব্যবহার করা যায়।

    \item আরও \textenglish{Full Adder} যুক্ত করে বড়
    \textenglish{bit-width}-এর \textenglish{Adder} তৈরি করা যায়।

\end{itemize}

\textbf{সীমাবদ্ধতা:}

একটি \textenglish{Full Adder}-এর \textenglish{Carry}
পরবর্তী \textenglish{Full Adder}-এ পৌঁছাতে কিছু সময় প্রয়োজন হয়।
ফলে \textenglish{stage} সংখ্যা বাড়লে
\textenglish{Carry propagation}-এর কারণে মোট
\textenglish{delay} বাড়ে।

\par\begingroup

\textbf{মনে রাখবে:}
\textenglish{Parallel Adder}-এ ইনপুট বিটগুলো একই সময়ে সংশ্লিষ্ট
\textenglish{Full Adder}-এ প্রবেশ করলেও
\textenglish{Carry} এক \textenglish{stage} থেকে পরবর্তী
\textenglish{stage}-এ ক্রমান্বয়ে যায়।
\par\endgroup



% =========================================================
% Applications
% =========================================================

\subsection{বাইনারি অ্যাডারের ব্যবহার}

\begin{enumerate}[leftmargin=*,itemsep=2pt]

    \item কম্পিউটারের
    \textenglish{ALU (Arithmetic Logic Unit)}-তে।

    \item বহু-বিট বাইনারি সংখ্যা যোগ করতে।

    \item \textenglish{Processor}-এর
    \textenglish{arithmetic operation}-এ।

    \item \textenglish{Digital Calculator}-এ।

    \item \textenglish{Address Calculation}-এ।

    \item বিভিন্ন
    \textenglish{digital arithmetic circuit} তৈরিতে।

\end{enumerate}

\subsection{অ্যাডারের সংখ্যা এবং প্রতিটি ধাপের ফল}
বাহ্যিক carry-in না থাকলে সবচেয়ে ডান বিটে একটি half adder এবং বাকি $n-1$ বিটে full adder দিয়ে $n$-বিট যোগ করা যায়। সব ধাপে full adder ব্যবহার করাও যায়; তখন প্রথমটির carry-in শূন্য। বাহ্যিক carry-in গ্রহণ করতে হলে প্রথম ধাপেও full adder প্রয়োজন।

\chThreeExample{১}{চার বিটের যোগে carry propagation}
$(1101)_2+(0111)_2$ এবং $C_0=0$।
\begin{center}\begin{tabular}{|c|c|c|c|c|c|}\hline
বিট $i$ & $A_i$ & $B_i$ & $C_i$ & $S_i$ & $C_{i+1}$\\\hline
0&1&1&0&0&1\\\hline
1&0&1&1&0&1\\\hline
2&1&1&1&1&1\\\hline
3&1&0&1&0&1\\\hline
\end{tabular}\end{center}
\chThreeAnswer ফল $C_4S_3S_2S_1S_0=10100_2=20_{10}$। চার বিটের sum output শুধু $0100$; $C_4$ আলাদা রাখলে সম্পূর্ণ চিহ্নবিহীন যোগফল পাওয়া যায়।

\chThreeExample{২}{ধারণাভিত্তিক MCQ}
বাহ্যিক carry-in ছাড়া দুটি আট বিট সংখ্যা যোগে half adder-সহ সরল নির্মাণ কোনটি? (ক) ৮ half adder, (খ) ১ half ও ৭ full adder, (গ) ৭ half ও ১ full adder, (ঘ) ১ full adder।
\chThreeAnswer (খ)। পরের বিটগুলোতে আগের carry গ্রহণের জন্য full adder লাগে। আট full adder-ও একটি বিকল্প নির্মাণ, কিন্তু প্রশ্নের নির্দিষ্ট half-adder-সহ নির্মাণটি (খ)।

% Mahbub teaching-type follow-up, 2026-10-14
\chThreeExample{৩}{চার বিটের যোগে Sum ও Carry আলাদা}
একটি চার-বিট adder-এ $A=1101$, $B=1011$, প্রথম carry $C_0=0$। চারটি sum bit ও শেষ carry লেখো।
\chThreeAnswer ডান থেকে $1+1+0=10$, $0+1+1=10$, $1+0+1=10$, $1+1+1=11$। তাই $S_3S_2S_1S_0=1000$, $C_4=1$। পূর্ণ unsigned ফল $C_4S_3S_2S_1S_0=11000_2=24$; শুধু sum port-এ দেখা মান ৮। পাঁচ-বিট যোগফল এবং চার-বিট sum output পৃথক। signed হিসেবে input দুটি $-3$ ও $-5$ হলে চার-বিট sum $1000=-8$ সঠিক; carry থাকার অর্থ signed overflow নয়।

\chThreeExample{৪}{Sum ও Carry থেকে সম্ভাব্য input}
একটি half adder-এ $S=1,C=0$; আর একটি full adder-এ $S=1,C_{out}=1$। সম্ভাব্য input নির্ণয় করো।
\chThreeAnswer half adder-এর যোগফল ১, তাই input জোড়া 01 অথবা 10। শুধু output থেকে কোন রেখাটি ১ ছিল তা নিশ্চিত নয়। full adder-এর পূর্ণ যোগফল $2C_{out}+S=3$; তিনটি এক-বিট input-এর যোগফল ৩ হতে হলে $A=B=C_{in}=1$। বিপরীত সমস্যায় প্রতিটি output combination-এর সঙ্গে কয়টি input মেলে তা সারণি থেকে পরীক্ষা করতে হয়।

